\documentclass{situationist-preprint}

\title{The Sublime After the Information Age}
\msid{AES-10-a}
\msfamily{AES-10}
\msdiscipline{Aesthetics}
\mstype{Principal manuscript}
\mskeywords{sublime; Kant; mathematical sublime; magnitude; comprehension; information; disenchantment}

\begin{document}
\maketitle

\begin{abstract}
It is often said that the sublime belongs to a past sensibility. Nature has been mapped and measured, its dangers are managed, and the metaphysics of the supersensible that sustained Kant's account is no longer widely held. We argue that the sublime remains possible. On a reconstruction of Kant's mathematical sublime that does without his account of reason's supersensible vocation, an experience is sublime when it presents a magnitude that the imagination cannot comprehend in a single intuition, while thought grasps it as a determinate whole, and the felt failure of the one is accompanied by pleasure in the reach of the other. None of these conditions requires ignorance of measure; on the contrary, the third requires measure. Quantification therefore does not end the sublime, and the magnitudes that the modern world makes available, including informational ones such as the size of a genome, of an archive or of a space of possible games, satisfy the conditions as well as the starry heavens did. We distinguish the sublime from mere bewilderment, and reply to the objections that such magnitudes are not presented to sense and that familiarity extinguishes the response.
\end{abstract}

\section{The problem}

Longinus applied the term \emph{hypsos} to elevated discourse. \citet{Burke1757} relocated the sublime in objects that produce astonishment through vastness, obscurity, power and infinity, and connected it with a delight in terror at a safe distance. \citet[\S\S 23--29]{Kant1790} distinguished the mathematical sublime, occasioned by magnitude, from the dynamical sublime, occasioned by might, and located the source of the feeling not in the object but in the mind: the failure of imagination to comprehend a magnitude reveals a faculty of reason that thinks the whole, and so a vocation that exceeds nature.

Doubts about the sublime's survival are of three kinds. The first is metaphysical: if the feeling depends on a supersensible vocation of reason, a naturalist cannot have it. \citet{Forsey2007} argues that no coherent theory of the sublime is possible without a transcendent commitment of this kind. The second is historical: the vast and dangerous nature of the eighteenth century has been surveyed, photographed and domesticated, and what replaced it, the technological sublime that \citet{Nye1994} traces through bridges, dams and skyscrapers, has itself become ordinary. The third concerns representation: today's magnitudes are known through numbers and images, not met in experience. Defenders of the concept have replied by distinguishing it from beauty and giving it a role in environmental aesthetics \citep{Brady2013,Shapshay2013}, or by interpreting Kant's account so that it reveals freedom without heavy metaphysics \citep{Clewis2009}.

We argue that the sublime remains available, specifically in the experience of magnitudes that exceed comprehension, and that informational magnitudes are among them.

\section{Framework}

\begin{definition}[Apprehension and comprehension]\label{d:comp}
To \emph{apprehend} a magnitude is to take up its parts successively in imagination; to \emph{comprehend} it is to hold the parts together as a single intuited whole.
\end{definition}

Kant's observation is that apprehension can proceed without limit, while comprehension reaches a maximum: as one adds parts, the first begin to fade before the last are reached \citep[\S 26]{Kant1790}.

\begin{definition}[The sublime]\label{d:sub}
An experience is \emph{sublime} if (i) it presents an object as having a magnitude; (ii) the subject's attempt to comprehend that magnitude fails; (iii) the subject nevertheless grasps the magnitude as a determinate whole in thought, by a number, a law or a concept of its structure; and (iv) the felt failure in (ii) is accompanied by pleasure in the grasp in (iii).
\end{definition}

Definition~\ref{d:sub} keeps Kant's structure and drops his explanation of (iv). For Kant the pleasure is a feeling of respect for our rational vocation. We require only that the subject take pleasure in the fact that thought reaches what imagination cannot; whether that fact reveals a supersensible faculty is a further question on which the experience does not depend.

\begin{principle}[No ignorance condition]\label{p:noign}
None of the conditions of Definition~\ref{d:sub} requires that the subject be ignorant of the measure of the object.
\end{principle}

\ref{p:noign} can be read off the definition: conditions (i), (ii) and (iv) say nothing about ignorance, and condition (iii) requires a grasp of the whole, which measure supplies.

\section{The argument}

\begin{proposition}[Measure enables the sublime]\label{p:measure}
Knowledge of the measure of a magnitude that exceeds comprehension makes a sublime experience of it possible rather than impossible.
\end{proposition}

\begin{proof}[Argument]
By \ref{p:noign}, ignorance of measure is not required, and condition (iii) requires that the whole be grasped in thought. A subject who does not know how large a thing is may be bewildered by it, but bewilderment lacks (iii): it is the failure of comprehension without the grasp that answers it. Kant's own culminating example confirms this. The progression from a tree to a mountain, to the earth's diameter, to the solar system, to the Milky Way and to systems of such galaxies is a progression of measures, and it produces the sublime by setting each comprehended unit against a whole that is grasped as a number of such units \citep[\S 26]{Kant1790}.
\end{proof}

Proposition~\ref{p:measure} answers the historical doubt. The surveying of nature replaces ignorance with measure, and measure is what the mathematical sublime requires. A subject who knows that the light now reaching her eye left the star before there were humans has more of what condition (iii) asks for, not less. This parallels the conclusion of AES-11, on which the appreciation of nature requires scientific rather than art-historical knowledge.

\begin{definition}[Informational magnitude]\label{d:info}
An \emph{informational magnitude} is the size of a body or space of information: the number of distinct items, configurations or possibilities it contains.
\end{definition}

\begin{theorem}\label{t:main}
Informational magnitudes can be objects of sublime experience.
\end{theorem}

\begin{proof}[Argument]
Consider three cases. The human genome contains about three billion base pairs. A subject who reads a stretch of it, letter by letter, apprehends its parts, and comprehension fails within a line; yet the whole is grasped as a sequence of that length, whose structure is in part understood. The number of possible games of chess was estimated by \citet{Shannon1950} at about $10^{120}$, a number exceeding the number of atoms in the observable universe; a player who knows the rules grasps the whole space by the rules that generate it, while no imagination comprehends even a small part of it. The Library of Babel imagined by \citet{Borges1941}, which contains every possible book of a fixed length, is a finite and determinate whole of which every reader holds the law, and which no reader could traverse. In each case the object is presented as having a magnitude, condition (i); comprehension fails, condition (ii); the whole is grasped as determinate by a number or a generating rule, condition (iii); and the subject who attends to the contrast can take pleasure in the reach of thought, condition (iv). The archives, networks and data sets of the present supply further instances of the same structure.
\end{proof}

\begin{corollary}\label{c:still}
The sublime remains possible.
\end{corollary}

\begin{proof}
Theorem~\ref{t:main} gives objects of the present for which all the conditions of Definition~\ref{d:sub} can be met, and Proposition~\ref{p:measure} shows that the measurement of nature has not removed the older ones.
\end{proof}

\section{Objections}

\begin{objection}[No presentation to sense]
For Kant the sublime requires that a magnitude be presented in intuition. An informational magnitude is known through a number, not seen; it is an object of calculation, not of aesthetic experience.
\end{objection}

\begin{reply}
Condition (ii) requires an attempt at comprehension, and such an attempt needs only apprehension, which can be sensory and successive: reading the sequence, scrolling through the archive, following a branching game tree move by move. Kant's own galaxies are not seen as galaxies either, but presented through a sensible unit set against a number. The informational cases differ only in that their units are symbols rather than spatial extents, and Kant counts magnitudes of any kind of homogeneous unit.
\end{reply}

\begin{objection}[Familiarity]
Large numbers have become routine. A subject who meets a new figure in the billions every day no longer feels anything at its scale.
\end{objection}

\begin{reply}
Habituation explains why the sublime is rare, which it always was: the eighteenth-century traveller who crossed the Alps more than once also ceased to tremble. The thesis is that the experience remains possible, and habituation is not an impossibility. It is undone by attention, by setting the figure against a comprehended unit, as Kant's progression does and as a reader of a genome is forced to do.
\end{reply}

\begin{objection}[Without the supersensible there is no sublime]
\citet{Forsey2007} argues that the sublime is defined by its relation to something beyond experience, and that a theory which drops this relation is a theory of something else.
\end{objection}

\begin{reply}
Definition~\ref{d:sub} retains the relation that matters for the experience: a whole beyond the reach of imagination, grasped by thought. What it drops is a claim about what that grasp reveals. The experience itself does not settle Kant's metaphysics; subjects who did not share it, from Burke's readers to modern astronomers, have reported the experience. A naturalist has the feeling, and only lacks one interpretation of it.
\end{reply}

\begin{objection}[Merely overwhelming]
The informational cases produce confusion and anxiety, not elevation. Their characteristic feeling is of being lost in data.
\end{objection}

\begin{reply}
Often they do, and then the experience is not sublime, because condition (iii) or (iv) fails. Being lost in data is the informational analogue of bewilderment before an unmeasured wilderness. The sublime requires that the whole be held in thought, by a law or a number, and it is available only to the subject who has that grasp.
\end{reply}

\section{Conclusion}

The sublime is the experience of a magnitude that imagination cannot comprehend and thought nevertheless grasps, with pleasure in the second answering displeasure in the first. Its conditions require measure rather than ignorance, so the measurement of nature has not ended it, and they are met by informational magnitudes as they are by astronomical ones. The sublime is still possible, and the present adds to its occasions.

\begin{thebibliography}{99}\small
\bibitem[Borges(1941)]{Borges1941} Borges, J.~L. (1941). La biblioteca de Babel. In \emph{El jard\'in de senderos que se bifurcan}. Sur.
\bibitem[Brady(2013)]{Brady2013} Brady, E. (2013). \emph{The Sublime in Modern Philosophy: Aesthetics, Ethics, and Nature}. Cambridge University Press.
\bibitem[Burke(1757)]{Burke1757} Burke, E. (1757). \emph{A Philosophical Enquiry into the Origin of Our Ideas of the Sublime and Beautiful}. Dodsley.
\bibitem[Clewis(2009)]{Clewis2009} Clewis, R.~R. (2009). \emph{The Kantian Sublime and the Revelation of Freedom}. Cambridge University Press.
\bibitem[Forsey(2007)]{Forsey2007} Forsey, J. (2007). Is a theory of the sublime possible? \emph{The Journal of Aesthetics and Art Criticism}, 65(4), 381--389.
\bibitem[Kant(1790)]{Kant1790} Kant, I. (1790). \emph{Kritik der Urteilskraft}. Lagarde und Friederich.
\bibitem[Nye(1994)]{Nye1994} Nye, D.~E. (1994). \emph{American Technological Sublime}. MIT Press.
\bibitem[Shannon(1950)]{Shannon1950} Shannon, C.~E. (1950). Programming a computer for playing chess. \emph{Philosophical Magazine}, 41(314), 256--275.
\bibitem[Shapshay(2013)]{Shapshay2013} Shapshay, S. (2013). Contemporary environmental aesthetics and the neglect of the sublime. \emph{British Journal of Aesthetics}, 53(2), 181--198.
\end{thebibliography}

\end{document}
